% ============================================================
%  KỊCH BẢN THUYẾT TRÌNH CHI TIẾT
%  Hệ Thống: XLA Face Privacy System
%  Nhóm: Đạt (Leader/Security) · Khánh (Backend/API) · Khoa (YOLO/Registry)
%  Thời lượng: ~30 phút
% ============================================================
\documentclass[12pt, a4paper]{article}

% --- Encoding & Language ---
\usepackage[utf8]{inputenc}
\usepackage[T5]{fontenc}
\usepackage[vietnamese]{babel}

% --- Layout ---
\usepackage[margin=2cm, top=2.5cm, bottom=2.5cm]{geometry}
\usepackage{parskip}
\usepackage{microtype}

% --- Color & Box ---
\usepackage[dvipsnames]{xcolor}
\usepackage{tcolorbox}
\tcbuselibrary{skins, breakable, listings}

% --- Code listing ---
\usepackage{listings}
\lstset{
  basicstyle=\ttfamily\small,
  keywordstyle=\color{NavyBlue}\bfseries,
  commentstyle=\color{OliveGreen}\itshape,
  stringstyle=\color{BrickRed},
  numbers=left,
  numberstyle=\tiny\color{gray},
  numbersep=6pt,
  frame=single,
  breaklines=true,
  showstringspaces=false,
  tabsize=4,
  backgroundcolor=\color{gray!8},
}

% --- Tables ---
\usepackage{booktabs}
\usepackage{tabularx}
\usepackage{array}
\usepackage{multirow}

% --- Icons / Symbols ---
\usepackage{amssymb}
\usepackage{pifont}
\usepackage{fontawesome5}
\newcommand{\cmark}{\ding{51}}
\newcommand{\xmark}{\ding{55}}

% --- Section styling ---
\usepackage{titlesec}
\titleformat{\section}[block]
  {\large\bfseries\color{NavyBlue}}
  {\thesection.}{0.5em}{}[\titlerule]
\titleformat{\subsection}[block]
  {\normalsize\bfseries\color{MidnightBlue}}
  {\thesubsection}{0.5em}{}

% --- Header / Footer ---
\usepackage{fancyhdr}
\pagestyle{fancy}
\fancyhf{}
\lhead{\textcolor{NavyBlue}{\textbf{XLA Face Privacy System}}}
\rhead{\textcolor{gray}{Kịch Bản Thuyết Trình}}
\cfoot{\thepage}

% --- Hyperlinks ---
\usepackage[hidelinks]{hyperref}

% --- Custom environments ---

% Dialog box (lời thoại của người trình bày)
\newtcolorbox{dialog}[2][]{
  title={\textbf{#2}},
  colback=blue!4,
  colframe=NavyBlue!60,
  fonttitle=\small\bfseries,
  left=4pt, right=4pt, top=4pt, bottom=4pt,
  breakable,
  #1
}

% Action box (hành động thao tác)
\newtcolorbox{action}[1][]{
  title={\faCogs\ \textbf{THAO TÁC}},
  colback=OliveGreen!6,
  colframe=OliveGreen!70,
  fonttitle=\small\bfseries,
  left=4pt, right=4pt,
  breakable,
  #1
}

% Note box (ghi chú kỹ thuật)
\newtcolorbox{note}[1][]{
  title={\textbf{LƯU Ý}},
  colback=Goldenrod!10,
  colframe=Goldenrod!80,
  fonttitle=\small\bfseries,
  left=4pt, right=4pt,
  breakable,
  #1
}

% Warning box
\newtcolorbox{warn}[1][]{
  title={\textbf{CẢNH BÁO}},
  colback=red!5,
  colframe=red!60,
  fonttitle=\small\bfseries,
  left=4pt, right=4pt,
  breakable,
  #1
}

% Timing badge macro
\newcommand{\timing}[1]{%
  \hfill\colorbox{NavyBlue!20}{\texttt{\small #1}}%
}

% Person label macros
\newcommand{\dat}{\textcolor{BrickRed}{\textbf{[ĐẠT]}}}
\newcommand{\khanh}{\textcolor{MidnightBlue}{\textbf{[KHÁNH]}}}
\newcommand{\khoa}{\textcolor{OliveGreen}{\textbf{[KHOA]}}}

% File path macro
\newcommand{\fpath}[1]{\texttt{\textcolor{purple}{#1}}}

% ============================================================
\begin{document}
% ============================================================

% ---- TITLE PAGE ----
\begin{titlepage}
\centering
\vspace*{2cm}
{\Huge\bfseries\color{NavyBlue} Kịch Bản Thuyết Trình Chi Tiết\par}
\vspace{0.5cm}
{\Large\color{MidnightBlue} XLA Face Privacy System\par}
\vspace{0.3cm}
{\large\color{gray} Môn Xử Lý Ảnh\par}
\vspace{1.5cm}

\begin{tcolorbox}[colback=NavyBlue!8, colframe=NavyBlue, width=0.85\textwidth]
\centering
\renewcommand{\arraystretch}{1.5}
\begin{tabular}{lll}
\textbf{Vai trò} & \textbf{Thành viên} & \textbf{Phụ trách} \\
\hline
\textcolor{BrickRed}{\textbf{Leader / Security}} & \textcolor{BrickRed}{\textbf{Đạt}} &
AES-256-GCM · PBKDF2 · Blur Engine \\
Backend / API & Khánh & FastAPI · Pipeline · Streaming \\
YOLO / Registry & Khoa & YOLOv11 · InsightFace · FaceDB \\
\end{tabular}
\end{tcolorbox}

\vspace{1.5cm}
\begin{tabular}{ll}
\textbf{Thời lượng:} & $\approx$ 30 phút \\
\textbf{Ngày:} & \today \\
\textbf{Repository:} & \texttt{github.com/Poicitaco/XLA\_T3} \\
\textbf{Branch:} & \texttt{feature/cipher} \\
\end{tabular}

\vspace{1.5cm}
\begin{note}
\textbf{Chuẩn bị trước khi trình bày:}
\begin{itemize}
  \item Server đang chạy: \texttt{uvicorn src.app\_api:app --host 0.0.0.0 --port 8000}
  \item Admin password đã được setup (POST /admin/setup)
  \item Có ít nhất 2 clip \texttt{.enc} trong \fpath{data/clips/secure/}
  \item VS Code mở sẵn: \fpath{src/clip\_recorder.py}, \fpath{src/admin\_auth.py}, \fpath{src/privacy.py}, \fpath{src/face\_recognition.py}
  \item Browser: 3 tab sẵn sàng tại \texttt{localhost:8000}
  \item Font size VS Code $\geq 16$, zoom browser 110\%
\end{itemize}
\end{note}
\end{titlepage}

% ---- TABLE OF CONTENTS ----
\tableofcontents
\newpage

% ============================================================
\section{Tổng Quan Kịch Bản \timing{0:00 – 30:00}}
% ============================================================

\begin{tcolorbox}[colback=NavyBlue!5, colframe=NavyBlue]
\centering
\renewcommand{\arraystretch}{1.6}
\begin{tabularx}{\textwidth}{clXl}
\toprule
\textbf{Phần} & \textbf{Người trình bày} & \textbf{Nội dung} & \textbf{Thời gian} \\
\midrule
1 & \textcolor{BrickRed}{\textbf{Đạt}} & Mở đầu, vấn đề, kiến trúc tổng thể & 0:00--6:00 \\
2 & \textcolor{OliveGreen}{\textbf{Khoa}} & YOLO Detection + InsightFace + Face Registry & 6:00--12:00 \\
3 & \textcolor{MidnightBlue}{\textbf{Khánh}} & FastAPI Backend + Live Pipeline + Streaming & 12:00--18:00 \\
4 & \textcolor{BrickRed}{\textbf{Đạt}} & Security Layer (Blur/Encrypt/Decrypt) & 18:00--25:00 \\
5 & \textbf{Cả nhóm} & Demo end-to-end + Q\&A & 25:00--30:00 \\
\bottomrule
\end{tabularx}
\end{tcolorbox}

% ============================================================
\section{Phần 1 — Đạt: Mở Đầu và Kiến Trúc \timing{0:00 – 6:00}}
% ============================================================

\subsection{Đặt vấn đề \timing{0:00 – 1:30}}

\begin{dialog}{\dat\ Lời mở đầu}
"Xin chào thầy cô và các bạn. Em là Đạt, đại diện nhóm chúng em
trình bày đồ án môn Xử Lý Ảnh với đề tài:
\textbf{Hệ thống camera thông minh bảo vệ quyền riêng tư khuôn mặt theo thời gian thực.}"

\vspace{0.3em}
"Câu hỏi đặt ra: Khi camera an ninh quay mọi người đi qua
--- khách, người giúp việc, hàng xóm ---
họ có đồng ý bị nhận dạng và lưu trữ khuôn mặt không?
Theo luật GDPR và quy định bảo vệ dữ liệu Việt Nam,
câu trả lời là \textbf{không}.
Nhưng nếu xóa hết camera thì ta lại mất bằng chứng khi có sự cố."

\vspace{0.3em}
"Hệ thống của nhóm em giải quyết đúng mâu thuẫn này:
\textbf{bảo vệ quyền riêng tư theo mặc định, đồng thời không mất đi bằng chứng.}"
\end{dialog}

\begin{note}
Nói chắc, đứng thẳng, nhìn vào mắt thầy cô khi nói câu định nghĩa cuối.
Đây là câu ``hook'' quan trọng nhất trong phần mở đầu.
\end{note}

\subsection{Giới thiệu hai tầng bảo vệ \timing{1:30 – 3:00}}

\begin{action}
Chuyển sang slide \textbf{``Dual Storage Architecture''}.
Dùng bút laser hoặc tay chỉ vào từng tầng khi giải thích.
\end{action}

\begin{dialog}{\dat\ Hai tầng bảo vệ}
"Hệ thống hoạt động theo nguyên lý \textbf{Dual Storage}.
Mỗi khi camera phát hiện khuôn mặt, hệ thống ghi lại \textbf{đồng thời hai bản}:

\vspace{0.5em}
\textbf{Tầng 1 --- Blurred Stream:}
Video khuôn mặt đã bị che giấu, lưu định dạng MP4 thông thường.
Bất kỳ ai trong hệ thống đều có thể xem --- không phân biệt được danh tính người trong clip.

\vspace{0.5em}
\textbf{Tầng 2 --- Encrypted Store:}
Video gốc hoàn toàn nguyên vẹn, được mã hóa bằng AES-256-GCM
và lưu dưới dạng file \texttt{.enc}.
\textit{Chỉ} người quản trị biết mật khẩu mới giải mã được.
Dùng khi cần điều tra pháp lý."
\end{dialog}

\subsection{Kiến trúc tổng thể \timing{3:00 – 5:00}}

\begin{action}
Mở VS Code, trỏ vào cây thư mục bên trái theo từng dòng.
Đọc to tên file và giải thích ngắn.
\end{action}

\begin{dialog}{\dat\ Kiến trúc 3 tầng}
"Kiến trúc hệ thống gồm 3 tầng:"

\vspace{0.5em}
"\textbf{Tầng Detection} --- \texttt{src/detector.py}, \texttt{src/face\_recognition.py}:
Khoa phụ trách. YOLOv11 phát hiện khuôn mặt, InsightFace trích xuất embedding
để nhận diện thành viên gia đình."

\vspace{0.5em}
"\textbf{Tầng Backend} --- \texttt{src/app\_api.py}, \texttt{src/backend\_api.py}:
Khánh phụ trách. FastAPI xử lý real-time pipeline,
phân luồng video đến đúng module, phục vụ giao diện frontend qua REST và WebSocket."

\vspace{0.5em}
"\textbf{Tầng Security} --- \texttt{src/privacy.py}, \texttt{src/clip\_recorder.py}, \texttt{src/admin\_auth.py}:
Em phụ trách. Blur engine xử lý 7 chế độ che giấu,
mã hóa AES-256-GCM với key derive từ PBKDF2-SHA256."
\end{dialog}

\begin{action}
Show nhanh cấu trúc thư mục trong terminal:
\begin{lstlisting}[language=bash]
Get-ChildItem src/ | Select-Object Name
\end{lstlisting}
\end{action}

\subsection{Chuyển sang Khoa \timing{5:30 – 6:00}}

\begin{dialog}{\dat\ Chuyển phần}
"Bây giờ em mời \textbf{Khoa} trình bày về phần nền tảng của pipeline ---
hệ thống phát hiện và nhận diện khuôn mặt.
Khoa ơi."
\end{dialog}

% ============================================================
\section{Phần 2 — Khoa: YOLO + InsightFace + Face Registry \timing{6:00 – 12:00}}
% ============================================================

\subsection{YOLOv11 Face Detection \timing{6:00 – 8:00}}

\begin{action}
Mở file \fpath{src/detector.py}. Scroll đến class \texttt{FaceDetector},
highlight phần \texttt{\_\_init\_\_} và method \texttt{detect\_faces}.
\end{action}

\begin{dialog}{\khoa\ YOLOv11 Detector}
"Em dùng \textbf{YOLOv11n-face} --- phiên bản nano của YOLO thế hệ 11,
được fine-tune đặc biệt để phát hiện khuôn mặt.
Model weights nằm tại \texttt{models/yolov11n-face.pt}."
\end{dialog}

\begin{action}
Chỉ vào file \fpath{src/detector.py}, highlight đoạn khởi tạo:
\end{action}

\begin{lstlisting}[language=Python, caption={src/detector.py --- Khởi tạo YOLO}]
class FaceDetector:
    def __init__(
        self,
        model_path: str = "models/yolov11n-face.pt",
        conf_threshold: float = 0.35,   # Nguong tin cay
        iou_threshold:  float = 0.50,   # Non-Maximum Suppression
        imgsz:          int   = 640,    # Kich thuoc anh inference
    ) -> None:
        yolo_cls = self._resolve_yolo_class()
        self.model = yolo_cls(model_path)
\end{lstlisting}

\begin{dialog}{\khoa\ Giải thích tham số}
"Ba tham số quan trọng:

\textbf{conf\_threshold = 0.35}: Chỉ chấp nhận detection có độ tin cậy trên 35\%.
Nếu để thấp quá sẽ có nhiều false positive --- nhầm đồ vật thành mặt người.

\textbf{iou\_threshold = 0.50}: Ngưỡng IoU cho Non-Maximum Suppression.
Khi hai bounding box chồng nhau quá 50\% thì giữ lại box có confidence cao hơn.

\textbf{imgsz = 640}: YOLO resize frame về 640×640 trước khi inference
để đảm bảo hiệu suất nhất quán bất kể độ phân giải camera."
\end{dialog}

\begin{action}
Scroll xuống method \texttt{detect\_faces}, highlight vòng lặp kết quả:
\end{action}

\begin{lstlisting}[language=Python, caption={src/detector.py --- detect\_faces}]
def detect_faces(self, frame: np.ndarray) -> List[List[int]]:
    results = self.model(
        frame,
        imgsz=self.imgsz,
        conf=self.conf_threshold,
        iou=self.iou_threshold,
        verbose=False,
    )
    boxes = []
    for r in results:
        for box in r.boxes:
            x1, y1, x2, y2 = map(int, box.xyxy[0])
            # Chuyen tu [x1,y1,x2,y2] sang [x,y,w,h]
            boxes.append([x1, y1, x2 - x1, y2 - y1])
    return boxes
\end{lstlisting}

\begin{dialog}{\khoa\ Format bounding box}
"YOLO trả về tọa độ góc trên trái và góc dưới phải \texttt{[x1,y1,x2,y2]}.
Em convert sang định dạng \texttt{[x, y, width, height]}
--- định dạng phổ biến hơn khi vẽ rectangle với OpenCV."
\end{dialog}

\subsection{InsightFace Embedding \timing{8:00 – 10:00}}

\begin{action}
Mở file \fpath{src/face\_recognition.py}. Highlight class \texttt{FaceRecognizer}.
\end{action}

\begin{lstlisting}[language=Python, caption={src/face\_recognition.py --- InsightFace setup}]
from insightface.app import FaceAnalysis

class FaceRecognizer:
    def __init__(
        self,
        model: str = "buffalo_l",       # Model buffalo_l: 512-dim embedding
        det_size: tuple = (256, 256),   # Kich thuoc detect noi bo InsightFace
    ) -> None:
        self.app = FaceAnalysis(
            name=model,
            providers=["CPUExecutionProvider"]  # Chay tren CPU
        )
        self.app.prepare(ctx_id=0, det_size=det_size)
\end{lstlisting}

\begin{dialog}{\khoa\ InsightFace Buffalo-L}
"InsightFace \textbf{buffalo\_l} là model nhận diện khuôn mặt state-of-the-art,
tạo ra vector embedding \textbf{512 chiều} --- gọi là ``vân tay số'' của khuôn mặt.

Mỗi người có một vector embedding đặc trưng.
Để nhận diện, em tính \textbf{cosine similarity} giữa embedding đang quan sát
và embedding đã đăng ký trong database:"
\end{dialog}

\begin{lstlisting}[language=Python, caption={src/face\_recognition.py --- Cosine Similarity}]
@staticmethod
def cosine_similarity(a: np.ndarray, b: np.ndarray) -> float:
    a_norm = np.linalg.norm(a)
    b_norm = np.linalg.norm(b)
    if a_norm == 0.0 or b_norm == 0.0:
        return -1.0
    # Gia tri tu -1 (nguoc hoan toan) den +1 (giong hoan toan)
    return float(np.dot(a, b) / (a_norm * b_norm))
\end{lstlisting}

\begin{dialog}{\khoa\ Ngưỡng nhận diện}
"Giá trị similarity từ \textbf{-1 đến +1}.
Ngưỡng mặc định là \textbf{0.35}: nếu similarity $> 0.35$,
đây là thành viên đã đăng ký $\rightarrow$ \textit{không blur}.
Nếu similarity $\leq 0.35$ $\rightarrow$ người lạ $\rightarrow$ \textit{bị blur}."
\end{dialog}

\begin{note}
Minh hoạ trực quan: ``Nếu trong database em có vector của Khoa,
và camera nhìn thấy Khoa, similarity sẽ khoảng 0.85.
Một người lạ chưa đăng ký sẽ cho similarity âm hoặc gần 0.''
\end{note}

\subsection{Face Registry Demo \timing{10:00 – 12:00}}

\begin{action}
Mở browser, chuyển đến tab \texttt{http://localhost:8000/admin/face-registry.html}.
Chỉ vào từng thành phần của UI.
\end{action}

\begin{dialog}{\khoa\ Demo Face Registry}
"Đây là trang đăng ký thành viên.
Quy trình đăng ký rất đơn giản:
nhập tên thành viên, upload ảnh khuôn mặt.
Backend tự động trích xuất embedding và lưu vào database."
\end{dialog}

\begin{action}
Mở VS Code, trỏ vào thư mục \fpath{data/members/}:
\begin{lstlisting}[language=bash]
Get-ChildItem data/members/ -Recurse |
  Where-Object { -not $_.PSIsContainer } |
  Group-Object DirectoryName |
  Select-Object Name, Count
\end{lstlisting}
\end{action}

\begin{dialog}{\khoa\ Giải thích dữ liệu}
"Hiện tại database có 2 thành viên: \textbf{Đạt} và \textbf{Khoa},
mỗi người \textbf{100 ảnh mẫu}.
Nhiều ảnh mẫu từ nhiều góc độ và điều kiện ánh sáng khác nhau
giúp embedding ổn định hơn và giảm thiểu nhận diện sai.

File JSON lưu embedding được persist tại
\texttt{data/family\_members.json}."
\end{dialog}

\begin{action}
Show nhanh file \fpath{data/family\_members.json},
highlight cấu trúc: \texttt{member\_id}, \texttt{name}, \texttt{threshold}, \texttt{embeddings} (danh sách vector).
\end{action}

\begin{dialog}{\khoa\ Chuyển phần}
"Em xin bàn giao cho Khánh trình bày về Backend API
và hệ thống pipeline thời gian thực."
\end{dialog}

% ============================================================
\section{Phần 3 — Khánh: Backend API + Live Pipeline \timing{12:00 – 18:00}}
% ============================================================

\subsection{FastAPI Architecture \timing{12:00 – 14:00}}

\begin{action}
Mở file \fpath{src/app\_api.py}. Scroll lên đầu file,
highlight phần docstring liệt kê tất cả endpoint.
\end{action}

\begin{dialog}{\khanh\ Giới thiệu FastAPI}
"Backend em xây dựng bằng \textbf{FastAPI} ---
Python async web framework, tự động generate OpenAPI documentation.

Toàn bộ API được chia thành nhóm chức năng rõ ràng:"
\end{dialog}

\begin{lstlisting}[language=Python, caption={src/app\_api.py --- Endpoint catalog}]
# Health / Info
GET  /health          # Kiem tra server song
GET  /info            # Version, capabilities, mode list
GET  /presets         # 4 preset: fast/balanced/strong/strict

# Xu ly anh tinh (stateless)
POST /anonymize       # Cha mat trong 1 anh upload
POST /lock            # Ma hoa AES-GCM tung vung mat
POST /unlock          # Giai ma, phuc hoi mat goc
POST /recognize       # Nhan dien tat ca mat trong anh

# Camera va Pipeline
GET  /cameras         # Liet ke camera kha dung
POST /pipeline/start  # Bat pipeline camera
DELETE /pipeline/stop # Tat pipeline
GET  /pipeline/status # Trang thai + thong ke live
PATCH /pipeline/config# Cap nhat config khong can restart

# Streaming
GET  /stream/mjpeg    # MJPEG stream (dung cho <img src="">)
GET  /stream/sse      # Server-Sent Events
WS   /stream/ws       # WebSocket (?format=binary|json)

# Clip management
GET  /clips           # Danh sach clip da luu
POST /clips/decrypt   # Giai ma clip (admin only)
DELETE /clips/{id}    # Xoa mot clip
\end{lstlisting}

\begin{dialog}{\khanh\ CORS và Security}
"Một điểm quan trọng trong cấu hình backend:"
\end{dialog}

\begin{lstlisting}[language=Python, caption={src/app\_api.py --- CORS middleware}]
app.add_middleware(
    CORSMiddleware,
    allow_origins=["*"],     # Dev mode; production nen gioi han
    allow_methods=["*"],
    allow_headers=["*"],
)
\end{lstlisting}

\subsection{Live Pipeline Architecture \timing{14:00 – 16:00}}

\begin{action}
Mở browser, chuyển đến \texttt{http://localhost:8000/admin/dashboard.html}.
Chỉ cho thầy cô xem layout tổng thể trước khi start pipeline.
\end{action}

\begin{dialog}{\khanh\ Giải thích pipeline}
"Khi nhấn Start, backend khởi động một background thread.
Luồng xử lý như sau:"
\end{dialog}

\begin{lstlisting}[language=Python, caption={Pipeline flow (minh họa)}]
# 1. Mo camera
cap = cv2.VideoCapture(camera_id)

# 2. FrameReader: doc frame tu camera theo thread rieng
# -> giam delay so voi doc dong bo

# 3. Moi frame:
if frame_index % detect_every == 0:
    boxes = detector.detect_faces(frame)   # YOLO
    boxes = smoother.smooth(boxes)         # EMA smoothing

# 4. RecognitionWorker: nhan dien async
#    So sanh embedding voi FamilyDatabase
#    -> identity: "Dat" | "Khoa" | None (nguoi la)

# 5. anonymize_faces():
#    Nhung mat co identity=None -> BLUR
#    Nhung mat co identity != None -> giu nguyen + ve ten

# 6. ClipRecorder.push_frame():
#    raw frame  -> .enc (AES-256-GCM)
#    display frame -> .mp4 (blurred)
\end{lstlisting}

\begin{action}
Thao tác live: chọn camera từ dropdown, mode \textbf{headcloak}, click \textbf{Start Pipeline}.
Chờ stream hiện lên, chỉ cho thầy cô thấy:
(a) người đã đăng ký không bị blur, (b) người lạ bị che hoàn toàn.
\end{action}

\begin{dialog}{\khanh\ Bbox Smoother}
"Một chi tiết kỹ thuật nhỏ nhưng quan trọng:
chúng em có module \texttt{bbox\_smoother.py} dùng \textbf{Exponential Moving Average}
để làm mịn bounding box qua các frame.

Không có smoother, bounding box sẽ nhảy lung tung frame-by-frame.
Với EMA alpha = 0.7, box di chuyển mượt mà và ổn định hơn nhiều."
\end{dialog}

\subsection{Hot-reload Config \timing{16:00 – 17:30}}

\begin{action}
Trong khi pipeline đang chạy, dùng UI đổi mode từ \texttt{headcloak} sang \texttt{pixelate}.
Chỉ cho thầy cô thấy thay đổi ngay lập tức trên stream.
\end{action}

\begin{lstlisting}[language=Python, caption={src/app\_api.py --- PATCH /pipeline/config}]
@app.patch("/pipeline/config")
async def update_config(req: ConfigUpdateRequest):
    """Cap nhat config pipeline KHONG can restart server."""
    global _pipeline_session
    if _pipeline_session is None:
        raise HTTPException(404, "Pipeline not running")

    # Atomic update: thay the config object in-place
    new_cfg = dc_replace(_pipeline_session.config, **req.dict(...))
    _pipeline_session.config = new_cfg
    return {"ok": True, "config": asdict(new_cfg)}
\end{lstlisting}

\begin{dialog}{\khanh\ Hot-reload}
"Endpoint PATCH \texttt{/pipeline/config} cho phép thay đổi
toàn bộ tham số của pipeline \textit{mà không cần tắt server}.
Đây là tính năng production-ready quan trọng ---
trong thực tế admin không thể restart camera mỗi lần cần đổi mode."
\end{dialog}

\subsection{Chuyển sang Đạt \timing{17:30 – 18:00}}

\begin{dialog}{\khanh\ Chuyển phần}
"Em xin bàn giao lại cho Đạt để trình bày phần quan trọng nhất của hệ thống ---
tầng bảo mật, mã hóa và decryption. Đạt ơi."
\end{dialog}

% ============================================================
\section{Phần 4 — Đạt: Security Layer \timing{18:00 – 25:00}}
% ============================================================

\begin{note}
Đây là phần \textbf{dài và quan trọng nhất} của toàn bài.
Đạt cần trình bày tự tin, chắc chắn, và giải thích đủ sâu
để thầy cô thấy rõ kiến thức bảo mật thực sự.
Đừng đọc code --- giải thích ý nghĩa từng dòng.
\end{note}

\subsection{Tại sao cần mã hóa? \timing{18:00 – 19:30}}

\begin{action}
Quay lại slide \textbf{``Dual Storage — Why Encryption?''}.
\end{action}

\begin{dialog}{\dat\ Động lực mã hóa}
"Khánh vừa demo pipeline blur rất đẹp.
Nhưng nếu chỉ dừng ở blur thì hệ thống có một lỗ hổng nghiêm trọng:
\textbf{bằng chứng bị mất}.

Giả sử tối nay có trộm đột nhập. Camera ghi lại clip.
Nhưng clip đã blur hoàn toàn --- khuôn mặt kẻ trộm không còn.
Video đó \textit{vô dụng} cho điều tra.

Giải pháp: lưu song song hai bản ---
bản public đã blur để mọi người xem bình thường,
và bản gốc mã hóa chỉ admin mới giải mã được.

Đây là triết lý thiết kế: \textbf{Privacy by Default, Evidence on Demand.}"
\end{dialog}

\subsection{AES-256-GCM Encryption \timing{19:30 – 21:30}}

\begin{action}
Mở file \fpath{src/clip\_recorder.py}. Scroll đến hàm \texttt{encrypt\_clip}
(khoảng dòng 110). Highlight toàn bộ hàm.
\end{action}

\begin{lstlisting}[language=Python, caption={src/clip\_recorder.py --- encrypt\_clip}]
def encrypt_clip(
    frames: List[np.ndarray],
    fps: float,
    password: str
) -> bytes:
    """
    Ma hoa frames bang AES-256-GCM.
    Output: [ 32B salt ][ 12B nonce ][ AES-GCM ciphertext + tag ]
    """
    # Buoc 1: Serialize frames thanh JPEG stream
    plaintext = _pack_frames(frames, fps)

    # Buoc 2: Sinh ngau nhien, moi clip mot salt RIENG
    salt  = secrets.token_bytes(32)   # 256-bit salt
    nonce = secrets.token_bytes(12)   # 96-bit nonce cho GCM

    # Buoc 3: Derive key tu password + salt qua PBKDF2
    key = derive_clip_key(password, salt)  # -> 32 bytes (AES-256)

    # Buoc 4: Ma hoa AES-256-GCM
    ct = AESGCM(key).encrypt(nonce, plaintext, None)

    # Buoc 5: Ghi file: salt | nonce | ciphertext
    return salt + nonce + ct
\end{lstlisting}

\begin{dialog}{\dat\ Giải thích từng bước}
"\textbf{Bước 1 --- Serialize:}
Frames được nén thành JPEG quality 90 rồi đóng gói
theo format length-prefixed: \texttt{[4B số frame][4B fps][4B len][jpeg bytes]...}
Tại sao JPEG 90? Đủ chất lượng để nhận diện pháp lý, file nhỏ hơn raw.

\vspace{0.5em}
\textbf{Bước 2 --- Random salt per clip:}
Điểm cực kỳ quan trọng: \textit{mỗi clip có một salt 32-byte ngẫu nhiên riêng}.
Dù admin dùng cùng một password,
key của clip A và clip B \textbf{hoàn toàn khác nhau}.
Hacker lấy được key của một clip cũng không dùng được cho clip khác.

\vspace{0.5em}
\textbf{Bước 3 --- PBKDF2 Key Derivation:}
Password $\rightarrow$ \texttt{PBKDF2-SHA256(password, salt, 390000 iterations)} $\rightarrow$ 32-byte key.
Tôi sẽ giải thích chi tiết phần này ngay sau.

\vspace{0.5em}
\textbf{Bước 4 --- AES-256-GCM:}
Tại sao \textit{GCM} chứ không phải CBC hay CTR?
GCM là \textbf{Authenticated Encryption} ---
ngoài mã hóa, còn tạo ra Authentication Tag để đảm bảo toàn vẹn.
Nếu ai sửa nội dung file \texttt{.enc} dù chỉ 1 bit,
decrypt sẽ fail ngay --- không trả về dữ liệu sai,
điều cực kỳ quan trọng với bằng chứng pháp lý.

\vspace{0.5em}
\textbf{Bước 5 --- File format:}
File \texttt{.enc} = \texttt{salt(32B)} + \texttt{nonce(12B)} + \texttt{ciphertext}.
Không có overhead gì --- compact và self-contained."
\end{dialog}

\subsection{PBKDF2-SHA256 Key Derivation \timing{21:30 – 23:00}}

\begin{action}
Mở file \fpath{src/admin\_auth.py}. Scroll đến đầu file,
highlight hằng số \texttt{\_ITERATIONS} và hàm \texttt{\_pbkdf2}.
\end{action}

\begin{lstlisting}[language=Python, caption={src/admin\_auth.py --- PBKDF2}]
# OWASP 2023 minimum for PBKDF2-SHA256
_ITERATIONS = 390_000
_KEY_LEN    = 32   # 256-bit AES key

def _pbkdf2(password: str, salt: bytes,
            iterations: int = _ITERATIONS) -> bytes:
    kdf = PBKDF2HMAC(
        algorithm=hashes.SHA256(),
        length=_KEY_LEN,      # Output: 32 bytes
        salt=salt,
        iterations=iterations,
    )
    return kdf.derive(password.encode("utf-8"))
\end{lstlisting}

\begin{dialog}{\dat\ Giải thích PBKDF2}
"\textbf{PBKDF2} là \textit{Password-Based Key Derivation Function 2.}
Nó biến một password ngắn, dễ nhớ thành một key mật mã học đủ mạnh.

Con số \textbf{390.000 iterations} không phải chọn ngẫu nhiên ---
đây là \textbf{chuẩn OWASP 2023} cho PBKDF2-SHA256.

Ý nghĩa thực tế:
Để thử một password, máy tính phải tính SHA256 \textit{390.000 lần}.
Với một máy hiện đại chạy 1 triệu hash/giây,
brute-force 8 ký tự mixed-case + số + symbol
tốn khoảng \textbf{hàng trăm ngàn năm}.

Và password admin \textbf{không bao giờ lưu dưới dạng plaintext}.
Nhìn vào file credentials:"
\end{dialog}

\begin{action}
Mở \fpath{data/admin\_credentials.json}, chỉ cho thầy cô xem.
\end{action}

\begin{lstlisting}[caption={data/admin\_credentials.json}]
{
  "salt": "base64-encoded-32-bytes...",
  "hash": "base64-encoded-pbkdf2-hash...",
  "iterations": 390000,
  "created_at": 1741234567.89
}
\end{lstlisting}

\begin{dialog}{\dat\ Credentials file}
"File này chỉ chứa \textbf{salt} và \textbf{hash}.
Không có password, không có key, không có gì có thể reverse.
Ngay cả em --- người viết code hệ thống ---
cũng \textit{không biết} password admin là gì
nếu không được admin cung cấp trực tiếp."
\end{dialog}

\subsection{Constant-Time Comparison \timing{23:00 – 23:45}}

\begin{action}
Trong file \fpath{src/admin\_auth.py}, highlight hàm \texttt{\_ct\_equal}.
\end{action}

\begin{lstlisting}[language=Python, caption={src/admin\_auth.py --- Timing Attack Prevention}]
def _ct_equal(a: bytes, b: bytes) -> bool:
    """
    Constant-time comparison - ngan Timing Attack.
    Python '==' dung lai ngay khi gap byte khac.
    hmac.compare_digest LUON so sanh het ca chuoi.
    """
    return hmac.compare_digest(a, b)
\end{lstlisting}

\begin{dialog}{\dat\ Timing Attack}
"Đây là một detail security nhỏ nhưng thể hiện tư duy bảo mật chuyên sâu.

Nếu dùng \texttt{a == b} bình thường, Python so sánh từng byte
và dừng lại ngay khi gặp byte khác nhau.
Hacker có thể đo thời gian response của server:
- Nếu response nhanh $\rightarrow$ sai từ byte đầu tiên
- Nếu response chậm hơn chút xíu $\rightarrow$ vài byte đầu đúng rồi

Qua hàng nghìn lần thử, hacker có thể đoán dần password đúng từng byte.
Đây gọi là \textbf{Timing Attack}.

\texttt{hmac.compare\_digest} luôn chạy một thời gian cố định
bất kể input khác nhau ở đâu $\rightarrow$ loại bỏ hoàn toàn timing leak."
\end{dialog}

\subsection{Blur Engine --- 7 Privacy Modes \timing{23:45 – 24:30}}

\begin{action}
Mở file \fpath{src/privacy.py}. Scroll nhanh qua
các hàm \texttt{\_fast\_gaussian\_blur}, \texttt{\_pixelate}, \texttt{\_head\_cloak},
\texttt{\_obliterate\_features} để thầy cô thấy sự đa dạng.
\end{action}

\begin{dialog}{\dat\ 7 Privacy Modes}
"Hệ thống hỗ trợ \textbf{7 chế độ ẩn danh}
từ nhẹ đến cực mạnh:"
\end{dialog}

\begin{tcolorbox}[colback=gray!5, colframe=NavyBlue!60, title=\textbf{7 Privacy Modes}]
\renewcommand{\arraystretch}{1.5}
\begin{tabularx}{\textwidth}{llXl}
\toprule
\textbf{Mode} & \textbf{Kỹ thuật} & \textbf{Mô tả} & \textbf{Tốc độ} \\
\midrule
\texttt{blur} & Gaussian Blur & Làm mờ Gaussian downscale 20\% & Nhanh \\
\texttt{pixelate} & Bilinear + Nearest & Pixel hóa mosaic 16px block & Nhanh \\
\texttt{solid} & Constant fill & Hình chữ nhật màu đặc \#282828 & Rất nhanh \\
\texttt{headcloak} & Mix + Noise + Blur & Che đầu + thêm noise ngẫu nhiên & Trung bình \\
\texttt{neckup} & Expand + Cloak & Che từ cổ trở lên & Trung bình \\
\texttt{obliterate} & 4-step pipeline & Compress $\rightarrow$ Scramble $\rightarrow$ Blur $\rightarrow$ Quantize & Chậm \\
\texttt{silhouette} & Extreme cloak & Randomize tone + noise + GBlur + Quantize & Chậm nhất \\
\bottomrule
\end{tabularx}
\end{tcolorbox}

\begin{lstlisting}[language=Python, caption={src/privacy.py --- obliterate (mạnh nhất)}]
def _obliterate_features(roi, obliterate_scale, ...):
    """4-buoc pha huy thong tin nhan dang - KHONG the khoi phuc"""

    # Buoc 1: Compress xuong 8% -> xoa het chi tiet mat
    work = cv2.resize(roi, (w*0.08, h*0.08), INTER_AREA)
    work = cv2.resize(work, (w, h), INTER_NEAREST)

    # Buoc 2: Xao tron block 12x12 -> pha vo quan he hinh hoc
    work = _scramble_blocks(work, block_size=12, seed=1337)

    # Buoc 3: Gaussian Blur 35x35 -> xoa canh con lai
    work = cv2.GaussianBlur(work, (35, 35), 0)

    # Buoc 4: Quantize mau ve 10 muc -> xoa texture tinh te
    work = _quantize_colors(work, levels=10)
    return work
\end{lstlisting}

\subsection{Demo Decrypt trên Secure Archive \timing{24:30 – 25:00}}

\begin{action}
Mở browser, chuyển sang tab \texttt{http://localhost:8000/admin/secure-archive.html}.
\end{action}

\begin{dialog}{\dat\ Demo decrypt}
"Đây là Secure Archive --- trang chỉ admin truy cập.
Danh sách các clip \texttt{.enc} đã được lưu.

Em sẽ decrypt một clip:"
\end{dialog}

\begin{action}
Thao tác:
\begin{enumerate}
  \item Click vào một clip trong danh sách
  \item Click \textbf{Decrypt \& Play}
  \item Nhập admin password vào modal
  \item Clip gốc (chưa blur) hiển thị
\end{enumerate}
\end{action}

\begin{dialog}{\dat\ Giải thích decrypt}
"Và đây là clip gốc --- khuôn mặt hoàn toàn rõ nét.

Nếu nhập sai password, server trả về lỗi:
\textit{`Decryption failed – wrong password or corrupted clip'}---
không leak bất kỳ thông tin nào về độ gần đúng của password.

Lưu ý: decrypt xảy ra \textit{server-side}.
File \texttt{.enc} không bao giờ gửi xuống browser.
Chỉ video đã decode được stream về --- thêm một lớp bảo vệ."
\end{dialog}

% ============================================================
\section{Phần 5 — Demo End-to-End + Q\&A \timing{25:00 – 30:00}}
% ============================================================

\subsection{Demo end-to-end hoàn chỉnh \timing{25:00 – 27:30}}

\begin{dialog}{\dat\ Giới thiệu demo}
"Bây giờ chúng em sẽ demo toàn bộ flow từ đầu đến cuối."
\end{dialog}

\subsubsection*{Bước 1: Reset và start pipeline}

\begin{action}
Khánh thao tác: stop pipeline nếu đang chạy, start lại với mode \textbf{obliterate}.
\end{action}

\begin{dialog}{\khanh\ }
"Pipeline đã start với chế độ obliterate --- mạnh nhất."
\end{dialog}

\subsubsection*{Bước 2: Thành viên đã đăng ký}

\begin{action}
Khoa đứng trước camera.
\end{action}

\begin{dialog}{\dat\ }
"Khoa đã được đăng ký trong database.
Hệ thống nhận ra Khoa với similarity cao ---
khuôn mặt \textbf{không bị blur}, và tên hiển thị bên cạnh."
\end{dialog}

\begin{action}
Đạt đứng trước camera.
\end{action}

\begin{dialog}{\dat\ }
"Đạt cũng đã đăng ký --- nhận ra, không blur."
\end{dialog}

\subsubsection*{Bước 3: Người lạ}

\begin{action}
Khánh (chưa đăng ký) đứng trước camera.
\end{action}

\begin{dialog}{\dat\ }
"Khánh chưa đăng ký vào database.
Hệ thống không nhận ra $\rightarrow$ khuôn mặt bị \textbf{obliterate hoàn toàn}.
Không ai xem stream này có thể biết Khánh trông như thế nào."
\end{dialog}

\subsubsection*{Bước 4: Mở Secure Archive}

\begin{action}
Đạt mở \texttt{http://localhost:8000/admin/secure-archive.html},
decrypt clip vừa ghi.
\end{action}

\begin{dialog}{\dat\ }
"Clip vừa ghi đã tự động lưu vào \texttt{data/clips/secure/}.
Sau khi decrypt, ta thấy khuôn mặt đầy đủ của Khánh trong clip gốc.

\textbf{Đây chính là mục tiêu:}
Khánh được bảo vệ quyền riêng tư khi xem video công khai,
nhưng bằng chứng vẫn được bảo toàn hoàn chỉnh cho admin."
\end{dialog}

\subsection{Privacy Modes Showcase \timing{27:30 – 29:00}}

\begin{action}
Đạt lần lượt switch qua từng mode trên dashboard,
chỉ cho thầy cô thấy kết quả trực tiếp trên stream.
\end{action}

\begin{tcolorbox}[colback=NavyBlue!4, colframe=NavyBlue, title=\textbf{Demo sequence}]
\begin{enumerate}
  \item \texttt{blur} --- \textit{``Cơ bản, vẫn còn thấy outline khuôn mặt mờ mờ''}
  \item \texttt{pixelate} --- \textit{``Mosaic, style camera truyền hình. Nhanh hơn blur''}
  \item \texttt{headcloak} (default) --- \textit{``Che toàn đầu + noise. Mặc định của hệ thống''}
  \item \texttt{obliterate} --- \textit{``4-bước pipeline, không còn trace nào. Dùng cho bảo vệ cao nhất''}
  \item \texttt{silhouette} --- \textit{``Cực mạnh: che đầu và vai, quantize màu còn 6 mức''}
\end{enumerate}
\end{tcolorbox}

\subsection{Kết luận \timing{29:00 – 30:00}}

\begin{dialog}{\dat\ Kết luận}
"Tóm lại, hệ thống chúng em đạt được \textbf{3 mục tiêu chính}:

\vspace{0.5em}
\textbf{1. Privacy bởi vì mặc định (Privacy by Default):}
Bất kỳ người lạ nào đi vào vùng quan sát đều bị ẩn danh hoá ngay lập tức,
không cần can thiệp thủ công.

\vspace{0.5em}
\textbf{2. Security đúng nghĩa:}
AES-256-GCM + PBKDF2-SHA256 390.000 iterations + constant-time comparison.
Không có backdoor, không có plaintext password, không có timing leak.

\vspace{0.5em}
\textbf{3. Usability thực tế:}
REST API đầy đủ cho frontend integration,
hot-reload config không cần restart,
7 chế độ blur để phù hợp nhiều use case.

\vspace{0.5em}
Nhóm chúng em xin cảm ơn thầy cô và \textbf{sẵn sàng trả lời câu hỏi}."
\end{dialog}

% ============================================================
\section{Câu Hỏi Thường Gặp --- Đáp Án Chuẩn Bị Sẵn}
% ============================================================

\begin{tcolorbox}[colback=gray!5, colframe=MidnightBlue, title=\textbf{Q1: Tại sao dùng AES-GCM mà không phải AES-CBC?},
  breakable]
\textbf{Trả lời (Đạt):}
AES-GCM là \textit{Authenticated Encryption with Associated Data (AEAD)}.
Ngoài mã hóa dữ liệu, GCM còn tạo ra một \textbf{Authentication Tag 128-bit}.
Khi decrypt, nếu ciphertext bị sửa đổi dù chỉ 1 bit, authentication sẽ fail.

AES-CBC chỉ mã hóa --- không đảm bảo toàn vẹn.
Kẻ tấn công có thể chỉnh sửa file \texttt{.enc}
và nhận lại dữ liệu sai mà không biết.

Với bằng chứng pháp lý, \textbf{integrity là bắt buộc, không phải tùy chọn}.
\end{tcolorbox}

\begin{tcolorbox}[colback=gray!5, colframe=MidnightBlue, title=\textbf{Q2: Nếu admin quên mật khẩu thì sao?},
  breakable]
\textbf{Trả lời (Đạt):}
Không khôi phục được --- đây là \textit{design có chủ ý}.

Hệ thống không lưu plaintext, không có secret recovery key, không có backdoor.
Đây là nguyên tắc \textbf{Zero Knowledge of Plaintext}.

Trong scope đồ án, ta chọn \textit{security over convenience}.
Trong production thực tế, có thể áp dụng thêm cơ chế
\textbf{Key Escrow} hoặc \textbf{Shamir Secret Sharing}
để chia key giữa nhiều người được ủy quyền.
\end{tcolorbox}

\begin{tcolorbox}[colback=gray!5, colframe=MidnightBlue, title=\textbf{Q3: Tại sao 390.000 iterations, không phải nhiều hơn?},
  breakable]
\textbf{Trả lời (Đạt):}
Con số này là khuyến nghị của \textbf{OWASP 2023} cho PBKDF2-SHA256
trên phần cứng phổ thông năm 2023.

Đây là \textit{balance}: đủ chậm để brute-force không khả thi,
nhưng đủ nhanh để user thực tế không cảm thấy lag khi encrypt/decrypt
một clip (xảy ra rất ít, không phải mỗi frame).

Nếu CPU server mạnh hơn, có thể tăng lên 600.000+ iterations
theo khuyến nghị mới hơn.
\end{tcolorbox}

\begin{tcolorbox}[colback=gray!5, colframe=MidnightBlue, title=\textbf{Q4: Real-time có bị lag không? FPS bao nhiêu?},
  breakable]
\textbf{Trả lời (Khánh):}
YOLO detect mỗi 2 frame (\texttt{detect\_every=2}).
Process width 640px với preset \texttt{balanced}.
Trên CPU mid-range đạt khoảng \textbf{15--20 FPS} --- đủ cho camera giám sát.

Nếu cần nhanh hơn:
Chuyển preset sang \texttt{fast}: YOLO size 416, pixelate, detect mỗi 3 frame,
có thể đạt 25--30 FPS.

Nếu có GPU, thay \texttt{CPUExecutionProvider} bằng \texttt{CUDAExecutionProvider}
cho InsightFace, tốc độ nhân diện bào nhận diện tăng gấp 5--10 lần.
\end{tcolorbox}

\begin{tcolorbox}[colback=gray!5, colframe=MidnightBlue, title=\textbf{Q5: Nhận diện nhầm thì sao? Privacy leakage?},
  breakable]
\textbf{Trả lời (Khoa):}
False negative (nhận diện sai thành viên là người lạ) $\rightarrow$
thành viên bị blur nhưng clip gốc vẫn có.
Không ảnh hưởng security.

False positive (người lạ nhầm thành thành viên) $\rightarrow$
người lạ không bị blur trong stream.
Đây là rủi ro privacy.

Để giảm thiểu, ngưỡng threshold được calibrate cẩn thận với 100 ảnh mẫu/người.
Threshold 0.35 được chọn sau benchmarking
để minimize false positive ở mức chấp nhận được.
\end{tcolorbox}

\begin{tcolorbox}[colback=gray!5, colframe=MidnightBlue, title=\textbf{Q6: GDPR / Luật bảo vệ dữ liệu?},
  breakable]
\textbf{Trả lời (Đạt):}
Hệ thống thiết kế theo nguyên tắc \textbf{Privacy by Design} (GDPR Article 25):
\begin{itemize}
  \item Dữ liệu biometric (embedding) chỉ lưu trong mạng nội bộ
  \item Clip gốc mã hóa, giới hạn truy cập
  \item Có thể delete clip và member bất cứ lúc nào
  \item Không gửi ảnh khuôn mặt ra bên ngoài server
\end{itemize}
Trong production cần bổ sung: audit log, consent management, data retention policy.
\end{tcolorbox}

% ============================================================
\section{Checklist Kỹ Thuật Trước Buổi Trình Bày}
% ============================================================

\begin{tcolorbox}[colback=OliveGreen!5, colframe=OliveGreen, title=\textbf{Checklist 30 phút trước}]
\begin{itemize}
  \item[$\square$] Server đang chạy trên port 8000
  \item[$\square$] Admin credentials tồn tại: \texttt{Test-Path data/admin\_credentials.json}
  \item[$\square$] Có ít nhất 2 clip \texttt{.enc} trong \texttt{data/clips/secure/}
  \item[$\square$] Camera nhận diện được: \texttt{GET http://localhost:8000/cameras}
  \item[$\square$] Thành viên Đạt và Khoa đã trong database
  \item[$\square$] VS Code: font size $\geq 16$, file explorer bên trái đóng bớt folder dư
  \item[$\square$] Browser: 3 tab mở sẵn, zoom 110\%
  \item[$\square$] Tắt notification Windows (Focus Assist: Alarms only)
  \item[$\square$] Tắt screensaver và sleep
  \item[$\square$] Backup: có sẵn clip \texttt{.enc} test nếu camera lỗi live
\end{itemize}
\end{tcolorbox}

\begin{tcolorbox}[colback=red!5, colframe=red!60, title=\textbf{Tình huống khẩn cấp}]
\textbf{Camera không hoạt động:}
Dùng clip demo quay sẵn. Giải thích: ``Do giới hạn phòng thi, em demo trước.''

\textbf{Server crash:}
\texttt{uvicorn src.app\_api:app --port 8000} --- khởi động lại $< 10$ giây.

\textbf{Quên từ:}
Không sao --- dừng, nhìn vào file code, giải thích từ code.
Code là ``notes'' của lập trình viên.

\textbf{Câu hỏi không biết:}
``Em chưa nghiên cứu sâu phần này,
nhưng theo kiến thức của em là... Thầy/cô có thể cho em bổ sung thêm không?''
\end{tcolorbox}

\vspace{1cm}
\begin{center}
\large\textbf{\textcolor{NavyBlue}{--- Chúc nhóm trình bày thành công! ---}}
\end{center}

\end{document}
